\section{Resumen y objetivos}
El Entregable 3 desarrolló una CNN para clasificar flores y analizó el conjunto
completo, muestras aleatorias reducidas y datos sintéticos obtenidos mediante
Distribution Matching. Su versión corregida incorporó también evidencia histórica
de hiperparámetros y regularización. El objetivo de esta entrega es ampliar esa
investigación y estudiar
su capacidad de ajuste y generalización. Se comparan una red totalmente conectada
(MLP), una CNN de cuatro convoluciones, una variante de seis convoluciones y
una CNN con AdamW y ajuste de tasa. Se cuantifican F1 macro, exactitud, pérdida,
parámetros y errores por clase. La MLP es una referencia reconstruida con el
pipeline del equipo para contrastar conectividad densa y convolucional: el
Entregable 3 corregido utilizó una CNN, no una MLP.

Para estudiar la regularización se incorpora, en un grupo experimental separado,
el mismo par de corridas históricas ya analizado en el A3 corregido, con y sin dropout, aumento de datos y decaimiento
de pesos. Se presentan tanto sus mejores checkpoints como una comparación en
la misma época registrada. Las curvas y las brechas permiten distinguir mayor
ajuste de train de una mejora efectiva en validación.
